\documentclass[12pt,letterpaper]{article}
\usepackage[utf8]{inputenc}
\usepackage[spanish]{babel}
\usepackage[margin=2.5cm]{geometry}
\usepackage{amsmath}
\usepackage{amsfonts}
\usepackage{amssymb}

\begin{document}

\begin{center}
    \textbf{\Large Examen de la Unidad II}\\
    \vspace{0.2cm}
    \textbf{Física para Ciencias de la Salud (005-1713)}\\
    \vspace{0.2cm}
    \textit{Cinemática, Dinámica, Trabajo, Energía y Potencia}
\end{center}

\vspace{0.5cm}
\noindent\textbf{Nombres y Apellidos:} \underline{\hspace{6cm}} \quad \textbf{C.I.:} \underline{\hspace{4cm}}
\vspace{0.5cm}

\noindent \textbf{Instrucciones:} Resuelva los siguientes 5 problemas de desarrollo utilizando una hoja en blanco para cada problema (2 puntos c/u). \textbf{Escriba de manera clara y que se entienda}. Muestre todo el procedimiento matemático, indique claramente las fórmulas utilizadas y exprese el resultado final con sus respectivas unidades del Sistema Internacional (S.I.). Considere $g = 9.8 \text{ m/s}^2$ cuando sea necesario.

\vspace{0.5cm}

\begin{enumerate}

    % === CINEMÁTICA ===
    \item \textbf{Cinemática:} Durante una sístole ventricular máxima, la sangre en la aorta pasa del reposo ($0 \text{ m/s}$) a una velocidad de $0.6 \text{ m/s}$ en un tiempo de $0.15 \text{ s}$. Calcule la aceleración media que experimenta el flujo sanguíneo.
    \vspace{0.5cm}

    % === DINÁMICA ===
    \item \textbf{Dinámica (Leyes de Newton):} En una sesión de rehabilitación fisioterapéutica, un paciente empuja un bloque de resistencia de $25 \text{ kg}$ sobre una superficie horizontal plana sin fricción. Si el paciente aplica una fuerza neta constante de $150 \text{ N}$, ¿qué aceleración adquiere el bloque?
    \vspace{0.5cm}

    % === TRABAJO MECÁNICO ===
    \item \textbf{Trabajo Mecánico:} Un paramédico ejerce una fuerza vertical hacia arriba de $400 \text{ N}$ a una velocidad constante para levantar una camilla desde el suelo hasta la altura de la ambulancia. Si el desplazamiento vertical de la camilla es de $0.8 \text{ m}$, calcule el trabajo mecánico realizado por el paramédico.
    \vspace{0.5cm}

    % === ENERGÍA POTENCIAL ===
    \item \textbf{Energía Potencial:} Se cuelga una bolsa de nutrición parenteral con una masa de $1.2 \text{ kg}$ en un atril, ubicándola a una altura de $1.5 \text{ m}$ por encima del punto de inserción en el brazo del paciente. Calcule la energía potencial gravitatoria de la bolsa respecto a dicho punto.
    \vspace{0.5cm}

    % === POTENCIA MECÁNICA ===
    \item \textbf{Potencia Mecánica:} Durante una prueba de esfuerzo ergoespirométrica, el corazón de un paciente realiza un trabajo mecánico de $75 \text{ Joules}$ en un lapso de $60 \text{ segundos}$. Determine la potencia media desarrollada por el corazón, expresada en Watts.
    \vspace{0.5cm}

\end{enumerate}

\end{document}